\documentclass[11pt]{amsart}

\usepackage{amsmath,amssymb,amsthm}
\usepackage{mathtools}
\usepackage[margin=1.1in]{geometry}
\usepackage[dvipsnames,svgnames]{xcolor}
\usepackage{hyperref}
\hypersetup{colorlinks=true,linkcolor=blue!55!black,citecolor=teal!60!black,urlcolor=blue!55!black,
  pdftitle={Emergent holonomy lives in degree zero},pdfauthor={MacBeth}}

\theoremstyle{plain}
\newtheorem{theorem}{Theorem}
\newtheorem{proposition}[theorem]{Proposition}
\newtheorem{lemma}[theorem]{Lemma}
\newtheorem{corollary}[theorem]{Corollary}
\theoremstyle{definition}
\newtheorem{definition}[theorem]{Definition}
\newtheorem{remark}[theorem]{Remark}
\newtheorem{example}[theorem]{Example}

\DeclareMathOperator{\Ind}{Ind}
\DeclareMathOperator{\Res}{Res}
\DeclareMathOperator{\Coind}{Coind}
\DeclareMathOperator{\Hom}{Hom}
\DeclareMathOperator{\Ext}{Ext}
\DeclareMathOperator{\Stab}{Stab}
\DeclareMathOperator{\Hh}{H}
\newcommand{\kG}{kG}
\newcommand{\kk}{k}
\newcommand{\conj}[2]{{}^{#1}\!#2}
\newcommand{\dcoset}[3]{#1\backslash #2/#3}

\title[Emergent holonomy lives in degree zero]{The emergent-holonomy count is a
degree-zero invariant:\\ $h(s)=\dim_k\Hom_{kU}(k[U/A],k[U/B])$}
\author{MacBeth}
\address{Kodamai}
\email{neil@kodamai.com}
\date{20 August 2026}

\begin{document}

\begin{abstract}
When two agents composed as a Zappa--Sz\'ep (exact) factorisation act on a shared
state and neither reserves the shared resource, an emergent holonomy appears: the two
factor-orbits through a state~$s$ cross in $h(s)$ points, and $h(s)=1$ exactly when the
composition is aligned. We identify this count with a representation-theoretic invariant
of the stabiliser group $U=\Stab_G(s)$. Writing $A,B\le U$ for the two factor-isotropies
and $k[U/A],k[U/B]$ for the associated permutation modules, the Mackey--Eckmann--Shapiro
decomposition of $\Ext$ between permutation modules gives
$h(s)=\dim_k\Hom_{kU}(k[U/A],k[U/B])=\dim_k\Ext^0_{kU}(k[U/A],k[U/B])$.
The novelty is not this classical formula but its consequence in the holonomy setting:
transversality of an exact factorisation forces every relevant subgroup intersection to
be trivial, so the entire higher $\Ext$-tower vanishes identically. The holonomy invariant
is therefore cohomological but concentrated in degree zero --- a single $\Hom$-space, or
equivalently a double-coset count --- with a theorem guaranteeing that no higher $\Ext$
obstruction can hide it. This corrects the working expectation that alignment should be
detected by a higher cohomology class, and explains why an orientation-twist of the
modules leaves the invariant unchanged.
We close by placing this vanishing inside a single formula. Over an \emph{abelian}
ambient group the Mackey decomposition factors as
$\Ext^n_{kG}(k[G/A],k[G/B])\cong\Hh^n(A\cap B;k)^{\oplus[G:AB]}$: a degree-independent
meeting count $h=[G:AB]$ times the cohomology of the overlap $A\cap B$. The holonomy
tower $[h,0,0,\dots]$ is the transverse corner $A\cap B=\{e\}$ of this formula; a
$V_4$ computation exhibits its full range, and support-variety geometry reads the
dichotomy as an intersection of rank varieties --- disjoint (transverse, vanishing)
versus coincident (non-transverse, surviving).
\end{abstract}

\maketitle

\section{Introduction}

\subsection*{Orchestration and emergent holonomy}
When two systems are composed, the order in which they act can matter. In the
compositional theory of directed containers, the sharpest instance of this is the
\emph{Zappa--Sz\'ep} (or \emph{exact}) factorisation of a group: $G=P\cdot P'$ with
$P\cap P'=\{e\}$ and $\lvert P\rvert\lvert P'\rvert=\lvert G\rvert$, so that every
$g\in G$ has a unique expression $g=p\,p'$ and a unique reverse expression $g=p''p'''$.
This is the algebraic skeleton of ``two agents sharing one state, each with its own
private moves,'' and it is the model we use throughout the container programme for
agent orchestration, supply chains, and stateful smart contracts; that composing two
directed containers this way is a Zappa--Sz\'ep product is due to
Ahman--Uustalu~\cite{au-distlaw}, and the emergent-holonomy content below sits on top of
their construction.

Fix a $G$-set $S$ and a state $s\in S$. The two factors $P,P'$ each sweep out an orbit
through $s$, namely $P\cdot s$ and $P'\cdot s$. Composing the agents means asking how
these two orbits interact, and the answer is a local invariant we call the
\emph{emergent holonomy}
\[
   h(s)\;:=\;\bigl\lvert (P\cdot s)\cap(P'\cdot s)\bigr\rvert
   \;=\;\lvert \dcoset{A}{U}{B}\rvert,
\]
the number of points where the two factor-orbits cross. Here $U=\Stab_G(s)$ is the
isotropy group of $s$ and $A=U\cap P$, $B=U\cap P'$ are the isotropies of the two
factors. That $h(s)$ is exactly the double-coset count $\lvert\dcoset{A}{U}{B}\rvert$,
and that $h(s)=1$ if and only if the composition is \emph{aligned} (the two orbits meet
only at $s$), was established in~\cite{macbeth-meeting}. Misalignment --- $h(s)>1$ ---
is the obstruction to treating the composite as a single well-defined move; it is the
container-theoretic signature of a reentrancy bug.

\subsection*{The question}
The count $h(s)$ is a combinatorial quantity. Is it \emph{cohomological}? Emergent
holonomy sits in a family of results where the obstruction to a clean composition is an
$\Hh^2$-class~\cite{macbeth-zs,macbeth-reentrancy}, and it is natural to expect the same
here: that alignment should be witnessed by the vanishing of some higher cohomology
class attached to the pair $(A,B)$ inside $U$. A concrete form of the expectation, due to
a collaborator, was that misalignment should surface as a nonzero class in
$\Ext^2_{kU}(k[U/A],k[U/B])$ or higher.

\subsection*{The result}
It does not. The holonomy invariant is representation-theoretic, but it lives entirely in
degree~zero, and the higher tower is not merely uninformative --- it is identically zero.

Let $k$ be a field, $M=k[U/A]$ and $N=k[U/B]$ the permutation modules on the coset spaces.
The classical Mackey--Eckmann--Shapiro decomposition of $\Ext$ between permutation modules
(Theorem~\ref{thm:mackey-ext}) reads
\[
   \Ext^n_{kU}(k[U/A],k[U/B])\;\cong\;
   \bigoplus_{u\in \dcoset{A}{U}{B}} \Hh^n\!\bigl(A\cap uBu^{-1};k\bigr),
   \qquad n\ge 0.
\]
In degree zero every summand is $\Hh^0(H;k)=k$, so the rank of $\Ext^0$ is the number of
double cosets. Combined with~\cite{macbeth-meeting} this is the dictionary:

\begin{theorem}[Dictionary, degree zero; Corollary~\ref{cor:dictionary}]\label{thm:intro-dict}
In the emergent-holonomy setting,
\[
  h(s)\;=\;\lvert\dcoset{A}{U}{B}\rvert
        \;=\;\dim_k\Hom_{kU}(k[U/A],k[U/B])
        \;=\;\dim_k\Ext^0_{kU}(k[U/A],k[U/B]).
\]
\end{theorem}

The point of this note is the next statement, which is what makes the degree-zero
location \emph{forced} rather than incidental.

\begin{theorem}[Degree-zero concentration; Theorem~\ref{thm:tower-vanishes}]\label{thm:intro-tower}
In the emergent-holonomy setting, for \emph{every} $u\in U$ the intersection
$A\cap uBu^{-1}$ is trivial. Consequently
\[
   \Ext^n_{kU}(k[U/A],k[U/B])\;\cong\;
   \begin{cases} k^{\,h(s)} & n=0,\\[2pt] 0 & n\ge 1,\end{cases}
\]
independently of $\operatorname{char}k$ and of whether the composition is aligned.
\end{theorem}

The mechanism is one line: $A\subseteq P$ and $uBu^{-1}\subseteq uP'u^{-1}$, so
$A\cap uBu^{-1}\subseteq P\cap uP'u^{-1}=\{e\}$ because an exact factorisation is
\emph{transverse} --- every conjugate of $P'$ meets $P$ trivially (Lemma~\ref{lem:transverse}).
Every Mackey summand in the decomposition is then the cohomology of the trivial group,
which is $k$ in degree $0$ and $0$ above.

\subsection*{What this corrects, and why it matters}
The working hypothesis --- ``a higher class detects alignment'' --- is false, and
Theorem~\ref{thm:intro-tower} says precisely why. The alignment content is the
\emph{rank of $\Ext^0$}: aligned $\iff h(s)=1\iff \Ext^0$ is one-dimensional. The higher
tower carries none of it. Our verification (Section~\ref{sec:verify}) includes a decisive
witness, computation~\textbf{W2}: an exact factorisation of $S_4$ with $h(s)=2>1$
(misaligned) whose entire higher $\Ext$-tower is nevertheless zero. A higher class cannot
be detecting what a misaligned example leaves invisible.

This also settles the collaborator's $\Ext^2$ conjecture in a clean way: it was right in
degree $0$ (where $h=\dim\Ext^0$) and structurally void in every higher degree, because
transversality leaves no higher $\Ext$ for a class to inhabit~\cite{macbeth-rickv4}. The
same transversality explains a twist-stability phenomenon (Corollary~\ref{cor:twist}):
tensoring $N$ with any character leaves the tower unchanged, so an ``orientation line''
cannot revive a higher obstruction.

For the wider programme the payoff is operational. Emergent holonomy of an unprotected
orchestration now carries a representation-theoretic signature,
$h(s)=\dim_k\Hom_{kU}(k[U/A],k[U/B])$, computable as a single space of equivariant maps ---
or, with no homological algebra at all, as a double-coset count --- \emph{together with a
guarantee} (Theorem~\ref{thm:intro-tower}) that no deeper $\Ext$ invariant can conceal the
obstruction. Detection is a degree-zero computation, and there is nothing further to audit.

\subsection*{On novelty}
The decomposition of Theorem~\ref{thm:mackey-ext} is classical: it is the Mackey formula
for $\Ext$ between permutation modules, going back to Nakaoka, and it is standard textbook
material~\cite{benson1,brown}. We give it in coordinates, with each ingredient cited,
because the container-theory audience for whom this note is written may not have the
group-cohomology backbone to hand. The contribution is the identification of $h(s)$ with
$\dim\Ext^0$ and the degree-zero concentration theorem, together with the corrected reading
they force.

\subsection*{Outline}
Section~\ref{sec:prelim} fixes notation for permutation modules, group cohomology, and the
exact-factorisation setup. Section~\ref{sec:mackey} assembles the three classical
ingredients into the Mackey $\Ext$-decomposition. Section~\ref{sec:dictionary} proves the
dictionary and the degree-zero concentration, and derives twist-stability.
Section~\ref{sec:verify} reports the independent numerical verification.
Section~\ref{sec:dichotomy} steps outside the holonomy locus to place the vanishing
inside a single abelian factorisation $\Ext=h\cdot\Hh^*(A\cap B)$, worked out on $V_4$
and read geometrically through support varieties.
Section~\ref{sec:conclusion} draws the operational conclusion and separates emergent
holonomy from the superficially similar higher towers that arise for unrelated subgroups.

\section{Preliminaries}\label{sec:prelim}

Throughout, $G$ is a finite group, $k$ a field of characteristic $p$ (the interesting case
is $p\mid\lvert G\rvert$), and all modules are \emph{left} modules over the group algebra
$kG$. For $g\in G$ and $B\le G$ write $\conj{g}{B}:=gBg^{-1}$.

\subsection{Induction, restriction, permutation modules}
For $H\le G$ and a $kH$-module $W$, induction and restriction are
\[
   \Ind_H^G W \;=\; kG\otimes_{kH} W,
   \qquad
   \Res_H Y \;=\; Y\ \text{regarded as a $kH$-module.}
\]
The \emph{permutation module} on the left cosets $G/H$ is
$k[G/H]=\Ind_H^G k$, the free $k$-module on $G/H$ with $G$ permuting the basis. We write
$M=k[G/A]=\Ind_A^G k$ and $N=k[G/B]=\Ind_B^G k$.

\subsection{Group cohomology}
For $H\le G$ we write $\Hh^n(H;k):=\Ext^n_{kH}(k,k)$ for group cohomology with trivial
coefficients. We use two standard facts. First, if $p\nmid\lvert H\rvert$ then $kH$ is
semisimple (Maschke), so $\Hh^{\ge 1}(H;k)=0$; only $p$-divisible subgroups contribute a
higher tower. Second, $\Hh^n(\{e\};k)=k$ for $n=0$ and $0$ for $n\ge 1$.

\subsection{Double cosets}
For $A,B\le G$, $\dcoset{A}{G}{B}$ denotes the set of double cosets $AgB$. A choice of
representatives $g$ is understood; different representatives of one double coset yield
conjugate intersection subgroups $A\cap\conj{g}{B}$, hence cohomologically isomorphic
summands, so every double-coset-indexed sum below is well defined.

\subsection{The emergent-holonomy setup}\label{ss:holonomy-setup}
An \emph{exact factorisation} (Zappa--Sz\'ep product) of $G$ is a pair of subgroups
$P,P'\le G$ with $P\cap P'=\{e\}$ and $PP'=G$; equivalently every $g\in G$ factors
uniquely as $g=pp'$ with $p\in P$, $p'\in P'$. Fix a $G$-set $S$ and $s\in S$, and set
\[
   U:=\Stab_G(s), \qquad A:=\Stab_P(s)=U\cap P, \qquad B:=\Stab_{P'}(s)=U\cap P'.
\]
Note $A\subseteq P$ and $B\subseteq P'$. The \emph{emergent holonomy} at $s$ is the
crossing count of the two factor-orbits,
\[
   h(s):=\bigl\lvert (P\cdot s)\cap(P'\cdot s)\bigr\rvert.
\]
We take as given the following identity, proved in~\cite{macbeth-meeting}.

\begin{proposition}[Meeting points~\cite{macbeth-meeting}]\label{prop:meeting}
$h(s)=\lvert\dcoset{A}{U}{B}\rvert=\lvert U\rvert/(\lvert A\rvert\,\lvert B\rvert)$, and
$h(s)=1$ if and only if the composition is aligned, i.e.\ $(P\cdot s)\cap(P'\cdot s)=\{s\}$.
\end{proposition}

The single structural fact about exact factorisations that we need is transversality.

\begin{lemma}[Transversality~\cite{macbeth-meeting}]\label{lem:transverse}
Let $G=P\cdot P'$ be an exact factorisation. Then $P\cap\conj{g}{P'}=\{e\}$ for every
$g\in G$.
\end{lemma}

\begin{proof}
Suppose $x\in P\cap gP'g^{-1}$, so $x=p\in P$ and $x=gp'g^{-1}$ for some $p'\in P'$.
Write $g=qq'$ with $q\in P$, $q'\in P'$ (unique factorisation). Conjugating $x$ by
$q^{-1}$ gives, on the one hand, $q^{-1}xq=q^{-1}pq\in P$, and on the other,
$q^{-1}xq=q^{-1}(qq')p'(qq')^{-1}q=q'p'q'^{-1}\in P'$. Hence
$q^{-1}xq\in P\cap P'=\{e\}$, so $x=e$.
\end{proof}

\begin{remark}
Lemma~\ref{lem:transverse} is the only place exactness is used. It says the factor $P$ is
transverse to \emph{every} conjugate of $P'$, not merely to $P'$ itself --- this global
transversality is exactly what will kill the higher $\Ext$-tower in Section~\ref{sec:dictionary}.
\end{remark}

\section{The Mackey decomposition of $\Ext$ between permutation modules}\label{sec:mackey}

This section assembles the classical decomposition. It rests on three ingredients, each a
standard lemma stated with a one-line justification; the references are
Benson~\cite[\S2.8,\ \S3.3]{benson1} and Brown~\cite[III.5--III.6]{brown}. Readers fluent in
group cohomology may skip to Theorem~\ref{thm:mackey-ext}.

\subsection{Eckmann--Shapiro}
\begin{lemma}[Eckmann--Shapiro]\label{lem:shapiro}
For $H\le G$, any $kH$-module $W$ and any $kG$-module $Y$,
\[
   \Ext^n_{kG}\!\bigl(\Ind_H^G W,\;Y\bigr)\;\cong\;\Ext^n_{kH}\!\bigl(W,\;\Res_H Y\bigr),
   \qquad n\ge 0.
\]
\end{lemma}
\begin{proof}
Choosing coset representatives $G=\bigsqcup_i g_i H$ exhibits $kG$ as a free right
$kH$-module, so $\Ind_H^G=kG\otimes_{kH}(-)$ is exact; and it sends free modules to free
modules ($\Ind_H^G kH=kG$), hence preserves projectives. Thus a projective resolution
$P_\bullet\to W$ over $kH$ induces a projective resolution $\Ind_H^G P_\bullet\to\Ind_H^G W$
over $kG$, and the adjunction $\Ind_H^G\dashv\Res_H$ gives a natural isomorphism of
complexes $\Hom_{kG}(\Ind_H^G P_\bullet,Y)\cong\Hom_{kH}(P_\bullet,\Res_H Y)$. Take
cohomology.
\end{proof}

Specialising $H=A$, $W=k$, $Y=N$ gives
\begin{equation}\label{eq:shapiro-spec}
   \Ext^n_{kG}(k[G/A],N)\;\cong\;\Ext^n_{kA}(k,\Res_A N)\;=\;\Hh^n(A;\Res_A N).
\end{equation}

\subsection{Mackey}
\begin{lemma}[Mackey decomposition]\label{lem:mackey}
As $kA$-modules,
\[
   \Res_A^G\Ind_B^G k\;\cong\;\bigoplus_{g\in \dcoset{A}{G}{B}}\Ind_{A\cap\conj{g}{B}}^A k .
\]
\end{lemma}
\begin{proof}
Decompose $G=\bigsqcup_{g} AgB$ into double cosets. As an $(A,B)$-set each double coset is
$AgB\cong A\times_{A\cap\conj{g}{B}}gB$, so the permutation $A$-set underlying $\Res_A(G/B)$
is $\bigsqcup_g A/(A\cap\conj{g}{B})$. Linearising gives the claim. (In general Mackey the
$g$-th summand carries a twist $\conj{g}{(-)}$ of the coefficient module; for the trivial
coefficient $k$ the twist is trivial.)
\end{proof}

Because $\Hh^n(A;-)$ commutes with finite direct sums in the coefficient module,
\eqref{eq:shapiro-spec} and Lemma~\ref{lem:mackey} combine to
\begin{equation}\label{eq:combined}
   \Ext^n_{kG}(k[G/A],k[G/B])\;\cong\;
   \bigoplus_{g\in\dcoset{A}{G}{B}}\Hh^n\!\bigl(A;\Ind_{A\cap\conj{g}{B}}^A k\bigr).
\end{equation}

\subsection{Shapiro for cohomology}
\begin{lemma}[Coinduction form]\label{lem:coind}
For $H\le A$, \; $\Hh^n(A;\Ind_H^A k)\cong \Hh^n(H;k)$.
\end{lemma}
\begin{proof}
For finite index, induction and coinduction coincide: $\Ind_H^A\cong\Coind_H^A$, both
$k[A/H]$ as a space. Now $\Coind_H^A$ is right adjoint to the exact functor $\Res_H$, which
preserves projectives ($kA$ is free over $kH$), so the derived adjunction gives
$\Ext^n_{kA}(k,\Coind_H^A k)\cong\Ext^n_{kH}(\Res_H k,k)=\Ext^n_{kH}(k,k)=\Hh^n(H;k)$.
\end{proof}

\subsection{The decomposition}
Substituting Lemma~\ref{lem:coind} (with $H=A\cap\conj{g}{B}$) into~\eqref{eq:combined}:

\begin{theorem}[Mackey--Eckmann--Shapiro; classical, cf.~{\cite[\S3]{benson1}}]\label{thm:mackey-ext}
For a finite group $G$, subgroups $A,B\le G$, and any field $k$,
\[
   \Ext^n_{kG}(k[G/A],k[G/B])\;\cong\;
   \bigoplus_{g\in\dcoset{A}{G}{B}}\Hh^n\!\bigl(A\cap gBg^{-1};k\bigr),
   \qquad n\ge 0,
\]
naturally in $n$, the sum over any set of double-coset representatives.
\end{theorem}

We claim no novelty for Theorem~\ref{thm:mackey-ext}; it is the Mackey formula for $\Ext$
of permutation modules (Nakaoka; see~\cite[\S3]{benson1}). Its degree-zero part is worth
recording separately, since it needs no resolution at all.

\begin{corollary}\label{cor:hom-doublecoset}
$\dim_k\Hom_{kG}(k[G/A],k[G/B])=\lvert\dcoset{A}{G}{B}\rvert$.
\end{corollary}
\begin{proof}
Take $n=0$ in Theorem~\ref{thm:mackey-ext}: every summand is $\Hh^0(-;k)=k$. Directly,
$\Hom_{kG}(k[G/A],N)=(\Res_A N)^A$, and the $A$-fixed points of the permutation module
$k[G/B]$ are spanned by the $A$-orbit sums on $G/B$, i.e.\ one basis vector per double
coset $AgB$.
\end{proof}

\section{The dictionary and degree-zero concentration}\label{sec:dictionary}

We now specialise to the emergent-holonomy setup of \S\ref{ss:holonomy-setup}, applying the
machinery of Section~\ref{sec:mackey} \emph{with $G$ replaced by the stabiliser group $U$}
and the subgroups $A=U\cap P$, $B=U\cap P'$.

\subsection{The dictionary}
\begin{corollary}[$h(s)=\dim\Ext^0$]\label{cor:dictionary}
In the emergent-holonomy setting,
\[
   h(s)\;=\;\lvert\dcoset{A}{U}{B}\rvert
       \;=\;\dim_k\Hom_{kU}(k[U/A],k[U/B])
       \;=\;\dim_k\Ext^0_{kU}(k[U/A],k[U/B]).
\]
\end{corollary}
\begin{proof}
The outer equality is Proposition~\ref{prop:meeting}; the inner two are
Corollary~\ref{cor:hom-doublecoset} applied over $U$ (with $\Ext^0=\Hom$).
\end{proof}

So the double-coset invariant one hoped $\Ext$ would see \emph{is} $\Ext$ --- in degree
zero. The remaining question is whether it is \emph{only} in degree zero.

\subsection{The higher tower vanishes}
\begin{theorem}[Degree-zero concentration]\label{thm:tower-vanishes}
In the emergent-holonomy setting, for every $u\in U$ the intersection $A\cap uBu^{-1}$ is
trivial. Consequently, over $U$,
\[
   \Ext^n_{kU}(k[U/A],k[U/B])\;\cong\;
   \bigoplus_{u\in\dcoset{A}{U}{B}}\Hh^n(\{e\};k)
   \;=\;
   \begin{cases} k^{\,h(s)} & n=0,\\[2pt] 0 & n\ge 1,\end{cases}
\]
independently of $\operatorname{char}k$ and of alignment.
\end{theorem}
\begin{proof}
Since $A\subseteq P$ and $uBu^{-1}\subseteq uP'u^{-1}$,
\[
   A\cap uBu^{-1}\;\subseteq\;P\cap uP'u^{-1}\;=\;\{e\}
\]
by Lemma~\ref{lem:transverse}, applicable because $u\in U\subseteq G$. Every Mackey
summand of Theorem~\ref{thm:mackey-ext} (over $U$) is therefore $\Hh^n(\{e\};k)$, which is
$k$ for $n=0$ and $0$ for $n\ge 1$. Summing over the $h(s)=\lvert\dcoset{A}{U}{B}\rvert$
double cosets gives the result.
\end{proof}

\begin{remark}[The corrected reading]\label{rem:reading}
Three consequences deserve to be stated plainly.
\begin{enumerate}
\item[(i)] The holonomy count is cohomological, but as a \emph{degree-zero rank}
   $h(s)=\dim\Ext^0$, not as a higher cohomology class. The tower is $[\,h,0,0,\dots\,]$;
   there is no $\Hh^2$-shadow of $h$ in it, and no higher class at all.
\item[(ii)] Alignment is read off in degree zero: aligned $\iff h(s)=1\iff\Ext^0$ is
   one-dimensional; misaligned $\iff h(s)>1\iff \dim\Ext^0>1$. Computation~\textbf{W2}
   below exhibits $h(s)=2>1$ (misaligned) with the higher tower still identically zero, so
   the higher tower demonstrably does \emph{not} detect alignment.
\item[(iii)] The transversality that forces the vanishing is exactly why an
   $\Ext^2$-shifted-class ansatz was structurally doomed~\cite{macbeth-rickv4}: in any
   exact-factorisation holonomy setting Lemma~\ref{lem:transverse} makes $\Res_A N$ free
   over $kA$, killing $\Ext^{\ge 1}$ --- there is no $\Ext^2$ for a class to live in. The
   ansatz was right in degree $0$ and void above it.
\end{enumerate}
\end{remark}

\subsection{Twist-stability}
The degree-zero concentration is robust under twisting the target module by a character,
which is the natural way one might try to inject an orientation into the invariant.

\begin{corollary}[Twist-stability]\label{cor:twist}
Let $L$ be any one-dimensional $kG$-module (character). Then
\[
   \Ext^n_{kG}(k[G/A],k[G/B]\otimes L)\;\cong\;
   \bigoplus_{g\in\dcoset{A}{G}{B}}\Hh^n\!\bigl(A\cap gBg^{-1};\,\Res_{A\cap gBg^{-1}}L\bigr).
\]
In the emergent-holonomy setting every $A\cap uBu^{-1}=\{e\}$, so $\Res L=k$ and the tower
is unchanged by the twist; over $\mathbb{F}_2$ there is no nontrivial character at all.
\end{corollary}
\begin{proof}
Mackey with coefficients gives $\Res_A(N\otimes L)=\bigoplus_g\Ind_{A\cap\conj{g}{B}}^A
\Res_{A\cap\conj{g}{B}}(\conj{g}{L})$; apply Lemma~\ref{lem:shapiro} and~\ref{lem:coind}.
Since $L$ is one-dimensional, $G$ acts on it by a character, which is conjugation-invariant,
so $\conj{g}{L}=L$. In the holonomy setting each intersection is trivial, so the restricted
coefficient is $k$.
\end{proof}

Thus an ``orientation line'' cannot revive a higher obstruction: the invariant is the rank
of $\Ext^0$ and nothing about a character changes it.

\section{Verification}\label{sec:verify}

Because Theorem~\ref{thm:mackey-ext} is assembled from several classical lemmas and the
whole point is the numerical value of the tower, we verified it independently rather than
trusting the assembly. The engine (\texttt{scratch/ext-mackey-general/}) implements its own
exact linear algebra over $\mathbb{F}_p$ --- no computer algebra system --- and computes
each side of Theorem~\ref{thm:mackey-ext} by a different route:
\begin{itemize}
\item \textbf{Left side:} a genuine minimal free resolution of $k[G/A]$ over $kG$, then
   $\Hh^n\Hom_{kG}(F_\bullet,k[G/B])$.
\item \textbf{Right side:} the Mackey double cosets $\dcoset{A}{G}{B}$, with each
   $\Hh^*(A\cap gBg^{-1};k)$ computed from its own $kH$-resolution.
\end{itemize}
All $17$ cases agree, and $\dim\Ext^0=\lvert\dcoset{A}{G}{B}\rvert$ in every case. The
direct left-side computation also pins the Mackey index as $A\cap gBg^{-1}$ (rather than
$gAg^{-1}\cap B$). A representative selection:

\begin{center}
\renewcommand{\arraystretch}{1.2}
\begin{tabular}{lccl}
\hline
Case & $p$ & $\Ext^{0,1,2,3}$ & note\\
\hline
$V_4,\ A=\langle a\rangle,B=\langle b\rangle$ & $2$ & $[1,0,0,0]$ & transverse, $1$ dcoset\\
$V_4,\ A=B=\langle a\rangle$ & $2$ & $[2,2,2,2]$ & $2$ dcosets, each $\cap=\mathbb{Z}/2$\\
$S_3,\ A=B=\langle(12)\rangle$ & $2$ & $[2,1,1,1]$ & mixed tower\\
$A_4,\ A=B=V_4$ & $2$ & $[3,6,9,12]$ & each $\cap=V_4$, $\Hh^n(V_4)=n{+}1$\\
$D_4,\ A=B=\langle\text{refl}\rangle$ & $2$ & $[3,2,2,2]$ & non-abelian $2$-group\\
$S_4,\ A=B=\langle(12)\rangle$ & $2$ & $[7,2,2,2]$ & $7$ dcosets\\
$\mathbb{Z}/3,\ A=B=\{e\}$ & $3$ & $[3,0,0,0]$ & regular, odd $p$\\
\hline
\textbf{W1}: $S_3{=}A_3{\cdot}\langle(12)\rangle$, $U{=}\langle(23)\rangle$ & $2$ & $[2,0,0,0]$ & $\Ext^0=h=2$\\
\textbf{W2}: $S_4{=}A_4{\cdot}\langle(12)\rangle$, $U{=}S_3,A{=}A_3,B{=}\{e\}$ & $2$ & $[2,0,0,0]$ & $h=2{>}1$, tower $\equiv 0$\\
\hline
\end{tabular}
\end{center}

The $A_4/V_4$ row is the strongest single stress-test of the assembly: the direct $kG$
resolution has Betti numbers $[1,2,3,4,5]$ and independently reproduces the
elementary-abelian tower $[3,6,9,12]$. The two rows below the rule are the holonomy
witnesses. \textbf{W1} confirms $\Ext^0=h=2$ with vanishing higher tower.
\textbf{W2} is the decisive one for Remark~\ref{rem:reading}(ii): it realises a misaligned
composition ($h=2>1$) whose higher $\Ext$-tower is identically zero --- a higher class
cannot be the alignment obstruction, because here alignment fails while every higher group
is zero.

\section{The full dichotomy: $\Ext=h\cdot\Hh^*(A\cap B)$ over an abelian group}\label{sec:dichotomy}

Theorem~\ref{thm:tower-vanishes} delivers the holonomy tower $[\,h,0,0,\dots\,]$ by
trivialising every Mackey intersection. It is worth seeing this vanishing not as an
isolated fact but as one corner of a single formula. When the ambient group is
\emph{abelian}, Theorem~\ref{thm:mackey-ext} collapses to a product of two independent
invariants --- a degree-independent multiplicity and a degree-graded shape --- and the
degree-zero concentration is exactly the case where the shape degenerates.

\subsection{The factorisation}
Let $G$ be a finite abelian group. Then $\conj{g}{B}=B$ for every $g$, so all the Mackey
summands of Theorem~\ref{thm:mackey-ext} share one intersection $A\cap B$; moreover
$AB\le G$ is a subgroup and each double coset $AgB=g(AB)$ is a coset of it, so
\[
   \lvert\dcoset{A}{G}{B}\rvert=[G:AB]=\frac{\lvert G\rvert\,\lvert A\cap B\rvert}
   {\lvert A\rvert\,\lvert B\rvert}.
\]
Writing $h:=[G:AB]$ for this meeting count, Theorem~\ref{thm:mackey-ext} reads:

\begin{corollary}[Abelian factorisation]\label{cor:abelian-factor}
For a finite abelian group $G$, subgroups $A,B\le G$, and any field $k$,
\[
   \Ext^n_{kG}(k[G/A],k[G/B])\;\cong\;\Hh^n(A\cap B;k)^{\oplus h},
   \qquad h=[G:AB],\quad n\ge 0,
\]
so that
\[
   \dim_k\Ext^n_{kG}(k[G/A],k[G/B])\;=\;
   \underbrace{[G:AB]}_{\text{meeting count }h}\;\cdot\;
   \underbrace{\dim_k\Hh^n(A\cap B;k)}_{\text{shape of the overlap}} .
\]
\end{corollary}
\begin{proof}
Abelian $G$ makes every $\conj{g}{B}=B$, so each of the $\lvert\dcoset{A}{G}{B}\rvert$
summands of Theorem~\ref{thm:mackey-ext} is $\Hh^n(A\cap B;k)$; the count is $[G:AB]$ as
computed above.
\end{proof}

The content is a separation of variables. The meeting count $h$ is a
\emph{degree-independent multiplicity}: it scales every degree of the tower equally, and
it is the same combinatorial number $\lvert\dcoset{A}{G}{B}\rvert$ that measures holonomy.
The cohomology $\Hh^*(A\cap B;k)$ is the \emph{per-degree shape}: it alone decides whether
a tower dies after degree zero, stays constant, or grows. Neither factor is new --- this is
Theorem~\ref{thm:mackey-ext} read off in the abelian case --- but the factored form is the
right way to see the dichotomy.

\subsection{Degree-zero concentration is the transverse corner}
The vanishing theorem is now the degenerate value of the shape factor.

\begin{corollary}[Transverse $=$ trivial overlap]\label{cor:transverse-corner}
In Corollary~\ref{cor:abelian-factor}, if $A\cap B=\{e\}$ then $\Hh^n(A\cap B;k)$ is $k$
for $n=0$ and $0$ for $n\ge 1$, so the tower is $[\,h,0,0,\dots\,]$. If instead
$A\cap B\ne\{e\}$ and $p\mid\lvert A\cap B\rvert$, the tower has a nonzero higher part.
\end{corollary}

For an abelian group an exact factorisation $G=A\cdot B$ is precisely a pair with
$A\cap B=\{e\}$ and $AB=G$, i.e.\ $h=1$; Corollary~\ref{cor:transverse-corner} then returns
$[1,0,0,\dots]$, the aligned instance of Theorem~\ref{thm:tower-vanishes}. So over an
abelian group the holonomy vanishing and the surviving non-transverse towers are two values
of one formula, switched by whether the overlap $A\cap B$ is trivial.

\begin{remark}
Theorem~\ref{thm:tower-vanishes} is \emph{not} subsumed by
Corollary~\ref{cor:abelian-factor}: for a non-abelian stabiliser $U$ the vanishing needs
$A\cap\conj{u}{B}=\{e\}$ for \emph{every} $u$, which Lemma~\ref{lem:transverse} supplies but
which a single-intersection statement does not. In the abelian examples below the two
coincide, and the abelian formula is what exhibits the surviving towers alongside the
vanishing one.
\end{remark}

\subsection{The worked example: $V_4$}
Take $G=V_4=\mathbb{Z}/2\times\mathbb{Z}/2=\langle a\rangle\times\langle b\rangle$ over
$k=\mathbb{F}_2$, so $kG=\mathbb{F}_2[x,y]/(x^2,y^2)$ with $x=a-1$, $y=b-1$ --- a local ring
with a unique simple module. Here $\dim\Hh^n(\mathbb{Z}/2;\mathbb{F}_2)=1$ in every degree
and $\dim\Hh^n(V_4;\mathbb{F}_2)=n+1$ (Poincar\'e series $(1-t)^{-2}$).
Corollary~\ref{cor:abelian-factor} predicts the four towers of Table~\ref{tab:v4}, and a
direct minimal-resolution computation confirms each in degrees $0$ through
$6$~\cite{macbeth-rickv4}.

\begin{table}[t]
\renewcommand{\arraystretch}{1.25}
\begin{tabular}{ccccll}
\hline
$A$ & $B$ & $A\cap B$ & $h$ & $\dim\Ext^n$, $n=0..6$ & regime\\
\hline
$\langle a\rangle$ & $\langle a\rangle$ & $\langle a\rangle$ & $2$ & $[2,2,2,2,2,2,2]$ & non-transverse, constant\\
$\langle a\rangle$ & $V_4$ & $\langle a\rangle$ & $1$ & $[1,1,1,1,1,1,1]$ & one $\Hh^*(\mathbb{Z}/2)$ copy\\
$\langle a\rangle$ & $\langle b\rangle$ & $\{e\}$ & $1$ & $[1,0,0,0,0,0,0]$ & transverse\\
$V_4$ & $V_4$ & $V_4$ & $1$ & $[1,2,3,4,5,6,7]$ & $=\Hh^*(V_4)$\\
\hline
\end{tabular}
\medskip
\caption{The four permutation-module Ext towers over $\mathbb{F}_2 V_4$, from
Corollary~\ref{cor:abelian-factor}. The multiplicity $h$ is a uniform vertical scaling:
rows one and two are the same shape $[1,1,1,\dots]$ at $h=2$ and $h=1$. The transverse row
is the aligned holonomy case of Theorem~\ref{thm:tower-vanishes}; the growing row is
$\Hh^*(V_4)$ itself.}\label{tab:v4}
\end{table}

The four rows span the whole range of the shape factor over a single group: a shape that
dies at degree zero (transverse), a shape that is constant ($\Hh^*(\mathbb{Z}/2)$, scaled by
$h\in\{1,2\}$), and a shape that grows linearly ($\Hh^*(V_4)$). Only the first is emergent
holonomy; the others are what the tower does when the overlap $A\cap B$ is allowed to be
nontrivial.

\subsection{Two structural readings}
The surviving tower raises two questions about \emph{where} the class lives --- one
representation-theoretic, one geometric. Both were posed by a collaborator; we record the
answers, the second invoking standard support-variety theory that we cite rather than
reprove.

\subsubsection*{Is the surviving class isolable to a single double coset?}
No: it is intrinsically a sum. In the constant case $A=B=\langle a\rangle$ there are $h=2$
double cosets, $\langle a\rangle\, e\,\langle a\rangle=\{e,a\}$ and
$\langle a\rangle\, b\,\langle a\rangle=\{b,ab\}$, and Corollary~\ref{cor:abelian-factor} is
a \emph{direct} sum in which \emph{each} coset carries an identical full copy of
$\Hh^*(\langle a\rangle;\mathbb{F}_2)=[1,1,1,\dots]$. The observed $\dim\Ext^n=2$ is $1+1$,
not one distinguished class of dimension $2$. The two summands are exchanged by the residual
symmetry $b\cdot(-)$, which swaps the cosets, so although the decomposition is canonical
(indexed by the double cosets) no single summand is preferred. When $h=1$ (rows two and
three of Table~\ref{tab:v4}) the tower does sit in the unique coset --- but that is the
degenerate case, not the general one.

\subsubsection*{On which stratum does it live?}
On a proper one. To a $kV_4$-module $M$ one attaches its \emph{rank variety}
$V_r(M)\subseteq\mathbb{P}^1$, the projectivised locus of $u=\alpha x+\beta y$ acting
non-freely on $M$; by the theory of Carlson and Avrunin--Scott the support of
$\Ext^*_{kG}(M,N)$ is the intersection $V_r(M)\cap
V_r(N)$~\cite{carlson-varieties,avrunin-scott,benson2}. For $M=k[V_4/\langle a\rangle]$
(dimension $2$) the generator $a$ acts trivially, so $x=a-1$ acts as the \emph{zero} matrix,
while $b$ swaps the two cosets, so $y=b-1$ has rank $1$; hence $u=\alpha x+\beta y=\beta(b-1)$
acts freely iff $\beta\ne 0$. Because $x$ acts as the literal zero matrix this is a
parametric identity over any field, giving
\[
   V_r\bigl(k[V_4/\langle a\rangle]\bigr)=\{(1{:}0)\}\subsetneq\mathbb{P}^1,
\]
a single point --- a proper, $0$-dimensional stratum. The dichotomy is then a variety
intersection:
\begin{itemize}
\item \textbf{transverse} ($A=\langle a\rangle$, $B=\langle b\rangle$): $V_r(M)=\{(1{:}0)\}$
   and $V_r(N)=\{(0{:}1)\}$ are \emph{disjoint}, so $\Ext^{\ge1}$ has empty support and
   vanishes --- the tower $[1,0,\dots]$ of Theorem~\ref{thm:tower-vanishes};
\item \textbf{non-transverse} ($A=B=\langle a\rangle$): both varieties are the \emph{same}
   point $\{(1{:}0)\}$, so the support is that point and $\Ext^{\ge1}\ne 0$.
\end{itemize}
The surviving class is thus not a top-stratum (complexity-$2$, full-support) class but a
complexity-$1$ class supported on the rank-$1$ stratum in the ``$a$-direction'' ---
consistent with its \emph{constant} dimension $\dim\Ext^n=2$ (polynomial growth of degree
$0$) and the periodic Betti numbers $[1,1,1,\dots]$ of the resolution of
$k[V_4/\langle a\rangle]$. Emergent holonomy, by contrast, lives where transversality forces
two such varieties apart: disjoint points, empty overlap, degree-zero tower. This is the
geometric shadow of Lemma~\ref{lem:transverse}.

\section{Conclusion}\label{sec:conclusion}

The emergent holonomy of an unprotected Zappa--Sz\'ep orchestration has a clean
representation-theoretic signature:
\[
   h(s)\;=\;\dim_k\Hom_{kU}\!\bigl(k[U/A],\,k[U/B]\bigr),
\]
the rank of the space of $U$-equivariant maps between the two factor-orbit permutation
modules. This is a degree-zero quantity, and Theorem~\ref{thm:tower-vanishes} guarantees
that it is the \emph{only} $\Ext$-invariant of the pair: transversality of the factorisation
forces the entire higher tower to vanish. Operationally, detecting misalignment is a single
$\Hom$-computation --- or, with no homological algebra, a double-coset count --- and one is
assured that no higher $\Ext$ obstruction lurks beneath it. There is nothing deeper to audit.

This corrects the natural but mistaken expectation, drawn from the $\Hh^2$-classified
obstructions elsewhere in the programme~\cite{macbeth-zs,macbeth-reentrancy}, that alignment
should be detected by a higher cohomology class. It is not. The alignment content is the
rank of $\Ext^0$, and the witness \textbf{W2} shows a misaligned composition with a
completely trivial higher tower. Consistently with the rigidity results
of~\cite{macbeth-twoomega}, there is no $\Hh^2$-class for $h$ in this tower.

\subsection*{A phenomenon not to conflate}
The higher $\Ext$-tower is not always zero; it vanishes here \emph{because} of the exact
factorisation. For \emph{general} subgroups $A,B\le G$ not arising from an exact
factorisation, Theorem~\ref{thm:mackey-ext} gives a genuine higher tower
$\bigoplus_g\Hh^{\ge 1}(A\cap gBg^{-1};k)$, supported on those double cosets with
$p\mid\lvert A\cap gBg^{-1}\rvert$; Section~\ref{sec:dichotomy} works this out in full for
abelian $G$, where it factors as $h\cdot\Hh^*(A\cap B)$ and emergent holonomy is the
transverse corner $A\cap B=\{e\}$. But the surviving tower is a different invariant ---
$p$-divisible \emph{overlap} between two subgroups --- and it must not be read as emergent
holonomy, which by Lemma~\ref{lem:transverse} always sits in the transverse locus where the
tower is zero.

\subsection*{Directions}
Two questions suggest themselves. First, the trivial higher tower says the pair
$(k[U/A],k[U/B])$ has $\Res_A N$ free over $kA$; it would be worth recording the resulting
$\Ext$-triviality as a structural statement about the container of factor-orbits, and asking
whether it characterises exact factorisations among all subgroup pairs with a given
$\Ext^0$. Second, the surviving tower $h\cdot\Hh^*(A\cap B)$ of
Section~\ref{sec:dichotomy} is a real invariant of \emph{overlap}, and its shape factor
$\Hh^*(A\cap B)$ is explicit for abelian $G$; whether this graded refinement of ``how badly
two unrelated agents share state'' has an orchestration meaning of its own --- and what the
shape factor becomes when $U$ is non-abelian, where $\conj{u}{B}$ genuinely varies --- is
open, and would sit alongside emergent holonomy rather than inside it.

\subsection*{Acknowledgements}
This note grew out of a collaboration on whether the holonomy count could be seen in
$\Ext^2$; the answer --- degree $0$, with a vanishing theorem above it --- is sharper than
either party first expected. Thanks to Neil Ghani for the container framing and to Robin
Langer for the infrastructure.

\begin{thebibliography}{9}

\bibitem{benson1}
D.~J. Benson,
\emph{Representations and Cohomology I: Basic Representation Theory of Finite Groups and
Associative Algebras},
Cambridge Studies in Advanced Mathematics 30, Cambridge University Press, 1991.
(Eckmann--Shapiro \S2.8; Mackey decomposition \S3.3; Mackey formula for $\Ext$ of
permutation modules \S3.)

\bibitem{brown}
K.~S. Brown,
\emph{Cohomology of Groups},
Graduate Texts in Mathematics 87, Springer, 1982. (Induction, restriction, and Shapiro's
lemma, III.5--III.6.)

\bibitem{benson2}
D.~J. Benson,
\emph{Representations and Cohomology II: Cohomology of Groups and Modules},
Cambridge Studies in Advanced Mathematics 31, Cambridge University Press, 1991.
(Rank and support varieties, Quillen stratification; $\operatorname{supp}\Ext^*(M,N)=
V_r(M)\cap V_r(N)$.)

\bibitem{carlson-varieties}
J.~F. Carlson,
\emph{The varieties and the cohomology ring of a module},
J. Algebra \textbf{85} (1983), 104--143.

\bibitem{avrunin-scott}
G.~S. Avrunin and L.~L. Scott,
\emph{Quillen stratification for modules},
Invent. Math. \textbf{66} (1982), 277--286.

\bibitem{macbeth-meeting}
MacBeth,
\emph{Emergent holonomy as meeting points: $h(s)=\lvert A\backslash U/B\rvert$},
Kodamai container programme, internal proof note, 13 August 2026.

\bibitem{au-distlaw}
D.~Ahman, T.~Uustalu,
\emph{Distributive laws of directed containers},
Progress in Informatics, No.~10 (2013), 3--18, DOI 10.2201/NIIPI.2013.10.2.
(Source of the Zappa--Sz\'ep-product construction: a distributive law of directed
containers assembles the joint container, generalizing the monoid Zappa--Sz\'ep product.)

\bibitem{macbeth-zs}
MacBeth,
\emph{Orchestration is a Zappa--Sz\'ep product \textnormal{(after Ahman--Uustalu
\cite{au-distlaw})}: the composition obstruction $[\omega]\in\Hh^2$},
Kodamai container programme, internal note, July 2026.

\bibitem{macbeth-reentrancy}
MacBeth,
\emph{Reentrancy and the class $[\omega]=\varepsilon$ in $\Hh^2\cong\mathbb{F}_2$}
(Lean-verified),
Kodamai container programme, internal note, July 2026.

\bibitem{macbeth-twoomega}
MacBeth,
\emph{Two $[\omega]$-sites are irreducibly distinct: skeletal rigidity kills the isotropy
restriction},
Kodamai container programme, internal proof note, 14 August 2026.

\bibitem{macbeth-rickv4}
MacBeth,
\emph{$\Ext^{\ge 1}$ vanishes for transverse subgroups: the $V_4$ twist-stable
computation},
Kodamai container programme, internal computation note, 20 August 2026.

\end{thebibliography}

\end{document}
